% GENERATED FILE. Edit canonical lessons, then run npm run generate:books.
% canonical-source-hash: fnv1a64:368a8a0d30153165
% canonical-lessons: SA-C480-dasati, SA-C480-visam, SA-C480-ksatam, SA-C480-asthibhangah, SA-C480-sravati

\chapter{Bites and Wounds}
\label{ch:sa-bites-and-wounds}
\hlchaptermodality{480}

\begin{chapteropening}
\textbf{By the end of this chapter:} \emph{I can say how someone was hurt.}

Complete the last lesson of chapter 480: I can say how someone was hurt.
\end{chapteropening}

\section[\texorpdfstring{\sk{दशति} (daśati)}{दशति (daśati)}]{\sk{दशति} (daśati) — bites}

\label{lesson:SA-C480-dasati}



\begin{quote}
Before the new one: say the Sanskrit for goes about, then the Sanskrit for a plague.
\end{quote}

\subsection*{You'll want to know: \sk{दशति}}
\textbf{\sk{दशति}} — \emph{daśati} — \textquotedblleft{}bites\textquotedblright{}.

\begin{quote}
\sk{सर्पः} \sk{पथिकं} \sk{दशति।}
\end{quote}

\begin{quote}
\emph{sarpaḥ pathikaṃ daśati.}
\end{quote}

\begin{quote}
The snake bites the traveller.
\end{quote}

Past (it …): \sk{शुनकः} \sk{बालं} \sk{दष्टवान्}, \emph{śunakaḥ bālaṃ daṣṭavān}.

\begin{cousinweb}
\textbf{\sk{दशति}} (\emph{daśati}) comes from the root √daṃś, \textquotedblleft{}to bite\textquotedblright{}, which drops its nasal in the present \emph{daśati} (MW).
\end{cousinweb}

\begin{grammarlens}[title={What you've built}]
One more word for describing what you lived and felt.
\end{grammarlens}

\subsection*{Your turn}
\begin{itemize}
\raggedright
  \item \emph{Say it:} \emph{daśati}
  \item \emph{Say it:} \emph{daśati}, once more
  \item \emph{Say it:} \emph{sarpaḥ pathikaṃ daśati.} — The snake bites the traveller.
  \item \emph{Recall:} say the Sanskrit for goes about, then the Sanskrit for a plague, then say \emph{daśati} again
\end{itemize}

\begin{tcolorbox}[breakable,colback=teal!4,colframe=teal!35!black,title={Before you move on}]
What does \sk{दशति} mean? (\textquotedblleft{}Bites\textquotedblright{}.) Say it once more.
\end{tcolorbox}
\section[\texorpdfstring{\sk{विषम्} (viṣam)}{विषम् (viṣam)}]{\sk{विषम्} (viṣam) — poison}

\label{lesson:SA-C480-visam}



\begin{quote}
Before the new one: say the Sanskrit for a plague, then the Sanskrit for bites.
\end{quote}

\subsection*{You'll want to know: \sk{विषम्}}
\textbf{\sk{विषम्}} — \emph{viṣam} — \textquotedblleft{}poison\textquotedblright{}.

\begin{quote}
\sk{वृश्चिकस्य} \sk{विषं} \sk{वेदनां} \sk{ददाति।}
\end{quote}

\begin{quote}
\emph{vṛścikasya viṣaṃ vedanāṃ dadāti.}
\end{quote}

\begin{quote}
A scorpion's poison causes agony.
\end{quote}

A neuter noun.

\begin{cousinweb}
\textbf{\sk{विषम्}} (\emph{viṣam}) MW glosses it \textquotedblleft{}anything active\textquotedblright{}, then \textquotedblleft{}poison, venom, bane\textquotedblright{}, and gives no root for it, so its origin is uncertain (MW).
\end{cousinweb}

\begin{grammarlens}[title={What you've built}]
One more word for describing what you lived and felt.
\end{grammarlens}

\subsection*{Your turn}
\begin{itemize}
\raggedright
  \item \emph{Say it:} \emph{viṣam}
  \item \emph{Say it:} \emph{viṣam}, once more
  \item \emph{Say it:} \emph{vṛścikasya viṣaṃ vedanāṃ dadāti.} — A scorpion's poison causes agony.
  \item \emph{Recall:} say the Sanskrit for a plague, then the Sanskrit for bites, then say \emph{viṣam} again
\end{itemize}

\begin{tcolorbox}[breakable,colback=teal!4,colframe=teal!35!black,title={Before you move on}]
What does \sk{विषम्} mean? (\textquotedblleft{}Poison\textquotedblright{}.) Say it once more.
\end{tcolorbox}
\section[\texorpdfstring{\sk{क्षतम्} (kṣatam)}{क्षतम् (kṣatam)}]{\sk{क्षतम्} (kṣatam) — an injury}

\label{lesson:SA-C480-ksatam}



\begin{quote}
Before the new one: say the Sanskrit for bites, then the Sanskrit for poison.
\end{quote}

\subsection*{You'll want to know: \sk{क्षतम्}}
\textbf{\sk{क्षतम्}} — \emph{kṣatam} — \textquotedblleft{}an injury\textquotedblright{}.

\begin{quote}
\sk{क्रीडायां} \sk{तस्य} \sk{पादे} \sk{क्षतम्} \sk{अभवत्।}
\end{quote}

\begin{quote}
\emph{krīḍāyāṃ tasya pāde kṣatam abhavat.}
\end{quote}

\begin{quote}
He hurt his foot while playing.
\end{quote}

A neuter noun.

\begin{cousinweb}
\textbf{\sk{क्षतम्}} (\emph{kṣatam}) comes from the root √kṣaṇ, \textquotedblleft{}to wound\textquotedblright{}; MW gives \textquotedblleft{}wounded, hurt, injured\textquotedblright{} and the noun \textquotedblleft{}a wound, an injury\textquotedblright{} (MW).
\end{cousinweb}

\begin{grammarlens}[title={What you've built}]
One more word for describing what you lived and felt.
\end{grammarlens}

\subsection*{Your turn}
\begin{itemize}
\raggedright
  \item \emph{Say it:} \emph{kṣatam}
  \item \emph{Say it:} \emph{kṣatam}, once more
  \item \emph{Say it:} \emph{krīḍāyāṃ tasya pāde kṣatam abhavat.} — He hurt his foot while playing.
  \item \emph{Recall:} say the Sanskrit for bites, then the Sanskrit for poison, then say \emph{kṣatam} again
\end{itemize}

\begin{tcolorbox}[breakable,colback=teal!4,colframe=teal!35!black,title={Before you move on}]
What does \sk{क्षतम्} mean? (\textquotedblleft{}An injury\textquotedblright{}.) Say it once more.
\end{tcolorbox}
\section[\texorpdfstring{\sk{अस्थिभंगः} (asthibhaṅgaḥ)}{अस्थिभंगः (asthibhaṅgaḥ)}]{\sk{अस्थिभंगः} (asthibhaṅgaḥ) — a broken bone}

\label{lesson:SA-C480-asthibhangah}



\begin{quote}
Before the new one: say the Sanskrit for poison, then the Sanskrit for an injury.
\end{quote}

\subsection*{You'll want to know: \sk{अस्थिभंगः}}
\textbf{\sk{अस्थिभंगः}} — \emph{asthibhaṅgaḥ} — \textquotedblleft{}a broken bone\textquotedblright{}.

\begin{quote}
\sk{वृद्धः} \sk{सोपानात्} \sk{अपतत्}, \sk{तस्य} \sk{अस्थिभंगः} \sk{अभवत्।}
\end{quote}

\begin{quote}
\emph{vṛddhaḥ sopānāt apatat, tasya asthibhaṅgaḥ abhavat.}
\end{quote}

\begin{quote}
The old man fell down the stairs and broke a bone.
\end{quote}

A masculine noun.

\begin{cousinweb}
\textbf{\sk{अस्थिभंगः}} (\emph{asthibhaṅgaḥ}) is built from \textbf{\sk{अस्थि}}, \textquotedblleft{}a bone\textquotedblright{}, and \emph{bhaṅga}, \textquotedblleft{}breaking\textquotedblright{}, from the root √bhañj; MW gives \textquotedblleft{}fracture of the bones\textquotedblright{} (MW).
\end{cousinweb}

\begin{grammarlens}[title={What you've built}]
One more word for describing what you lived and felt.
\end{grammarlens}

\subsection*{Your turn}
\begin{itemize}
\raggedright
  \item \emph{Say it:} \emph{asthibhaṅgaḥ}
  \item \emph{Say it:} \emph{asthibhaṅgaḥ}, once more
  \item \emph{Say it:} \emph{vṛddhaḥ sopānāt apatat, tasya asthibhaṅgaḥ abhavat.} — The old man fell down the stairs and broke a bone.
  \item \emph{Recall:} say the Sanskrit for poison, then the Sanskrit for an injury, then say \emph{asthibhaṅgaḥ} again
\end{itemize}

\begin{tcolorbox}[breakable,colback=teal!4,colframe=teal!35!black,title={Before you move on}]
What does \sk{अस्थिभंगः} mean? (\textquotedblleft{}A broken bone\textquotedblright{}.) Say it once more.
\end{tcolorbox}
\section[\texorpdfstring{\sk{स्रवति} (sravati)}{स्रवति (sravati)}]{\sk{स्रवति} (sravati) — oozes}

\label{lesson:SA-C480-sravati}



\begin{quote}
Before the new one: say the Sanskrit for an injury, then the Sanskrit for a broken bone.
\end{quote}

\subsection*{You'll want to know: \sk{स्रवति}}
\textbf{\sk{स्रवति}} — \emph{sravati} — \textquotedblleft{}oozes\textquotedblright{}.

\begin{quote}
\sk{व्रणात्} \sk{रक्तं} \sk{स्रवति।}
\end{quote}

\begin{quote}
\emph{vraṇāt raktaṃ sravati.}
\end{quote}

\begin{quote}
Blood oozes from the wound.
\end{quote}

Past (it …): \sk{रक्तं} \sk{व्रणात्} \sk{स्रुतवत्}, \emph{raktaṃ vraṇāt srutavat}.

\begin{cousinweb}
\textbf{\sk{स्रवति}} (\emph{sravati}) comes from the root √sru, \textquotedblleft{}to flow, to stream\textquotedblright{}; MW's ppp \emph{sruta} is \textquotedblleft{}streaming, flowing\textquotedblright{} (MW).
\end{cousinweb}

\begin{grammarlens}[title={What you've built}]
One more word for describing what you lived and felt.
\end{grammarlens}

\subsection*{Your turn}
\begin{itemize}
\raggedright
  \item \emph{Say it:} \emph{sravati}
  \item \emph{Say it:} \emph{sravati}, once more
  \item \emph{Say it:} \emph{vraṇāt raktaṃ sravati.} — Blood oozes from the wound.
  \item \emph{Recall:} say the Sanskrit for an injury, then the Sanskrit for a broken bone, then say \emph{sravati} again
\end{itemize}

\begin{tcolorbox}[breakable,colback=teal!4,colframe=teal!35!black,title={Before you move on}]
What does \sk{स्रवति} mean? (\textquotedblleft{}Oozes\textquotedblright{}.) Say it once more.
\end{tcolorbox}
